\begin{statement}{P3K.57}{Proved}
For the marked source
\[
 y^2=x^3+(-2+3\epsilon^2)x^2+t x
       +\epsilon^2(1-\epsilon^2)^2,\qquad
 G=(0,\epsilon(1-\epsilon^2)),
\]
retain the selected whole-zero fibre at
\(t_*=(1-\epsilon^2)(1-3\epsilon^2)\), the actual Cartier line
\(M_a=\mathcal O([aG]-[O])\), and the suspension over \(-bG\),
where \(ab=2\). Use either the initial vertical-edge cusp fan or its
literal reflected fan, and the admitted cap data
\[
 \begin{array}{c|c|c}
 (a,b)&\text{cap order}&\text{root type}\\ \hline
 (1,2)&m=4p&\text{good III-star root}\\
 (2,1)&m=4p-1&\text{reduced III root}
 \end{array}\qquad(p\ge1),
\]
with either norm sign \(n=\pm1\) and every retained integral
\(d\delta\) twist. These are the completed sphere families of
P3F.41 and P3K.24, with their prescribed source, line and cap maps.
For a fixed sufficiently small nonzero suspension scalar, the actual
positive \(\epsilon\)-loop has Jin characteristic integer and oriented
clutched-sphere class
\[
 \begin{array}{c|c|c}
 (a,b)&\displaystyle\int\frac{4x^3-p_1x}{24}
       &\Theta_7\simeq\mathbb Z/28\\ \hline
 (1,2)&0&0\\
 (2,1)&1&1.
 \end{array}
\]
The actual source/Cartier/cap morphism under
\(v=\epsilon^{-2}-1\) preserves the suspension scalar and retains
the full \(D^{-4}\) root endpoint, the reduced doubled lattice period,
the annular reference overlap and the quadratic scalar cochain.
The theorem assigns no source-zero return and makes no further fan comparison.
\end{statement}
\begin{proof}
\par\medskip\noindent\textbf{IF.1. The marked infinity chart and selected-node units.}\par\nopagebreak[4]

On the cover \(\rho^2=1+v\), put \(\epsilon=\rho^{-1}\). The source map is
\[
 v=\epsilon^{-2}-1,\qquad
 X=\epsilon^{-2}x,\quad Y=\epsilon^{-3}y,\quad T=\epsilon^{-4}t.
\]
It carries the actual marked equation and section to
\[
 y^2=x^3+(-2+3\epsilon^2)x^2+t x
                     +\epsilon^2(1-\epsilon^2)^2,\qquad
 G=(0,g),\quad g=\epsilon(1-\epsilon^2).                 \tag{P3K.57.IF.1}
\]
Write \(A=1-\epsilon^2\), \(t_*=A(1-3\epsilon^2)\), and \(q_*=t-t_*\).
The selected fibre factors and normalizes by the literal formulas
\[
 (x-A)^2(x+\epsilon^2),\qquad
 x=s^2-\epsilon^2,\quad y=s(s^2-1).                     \tag{P3K.57.IF.2}
\]
Its double preimages are \(s=\pm1\), independent of the parameter. The marked
section has \(s=-\epsilon\). The normalization derivative has first coordinate
\(2s\); at \(s=0\) its second coordinate is \(-1\). It is therefore an immersion
through the full parameter disc, with its two double sections disjoint.

Choose the nodal multiplicative coordinate \(Z=(s-1)/(s+1)\). Its marked point
is
\[
 P_G=-\frac{1+\epsilon}{1-\epsilon},\qquad
 P_G(-\epsilon)=P_G(\epsilon)^{-1}.                    \tag{P3K.57.IF.3}
\]
The selected Tate parameter satisfies
\[
 \left.\frac{Q}{q_*}\right|_{q_*=0}
       =-\frac{1-\epsilon^2}{64}.                       \tag{P3K.57.IF.4}
\]
Indeed the cubic discriminant has \(t\)-derivative \(-4A\) at \(t_*\), its
elliptic discriminant is sixteen times that cubic discriminant, and \(c_4=16\)
there. The first term of \(j(Q)^{-1}\) is \(Q\), which gives (P3K.57.IF.4).

Retain the actual two chord factors
\[
 F_{21}=\frac{y+g+t x/(2g)}x,\qquad
 F_{12}=\frac{y-g-t x/(2g)}x .
\]
At the selected node they are \(A/(2\epsilon)\) and
\(-A/(2\epsilon)\). The transported actual Hom map is
\[
 \phi_{ab}=\frac{q_*}{g F_{ab}},\qquad
 \left.\frac{\phi_{21}}{q_*}\right|_{\rm node}=\frac2{A^2},
 \quad
 \left.\frac{\phi_{12}}{q_*}\right|_{\rm node}=-\frac2{A^2}.       \tag{P3K.57.IF.5}
\]
All three units (P3K.57.IF.3)--\allowbreak (P3K.57.IF.5) extend and are nonzero on a sufficiently small
full \(\epsilon\)-disc. Their full-disc parameter windings are zero. This is
an actual selected-cusp chart statement; it does not trivialize the exterior
suspension at the other singular fibres.

\par\medskip\noindent\textbf{IF.2. The complete suspension keeps the scalar.}\par\nopagebreak[4]

Use the canonical meromorphic divisor sections to identify the old Cartier
line with the line of (P3K.57.IF.1). Translation by \(G\) commutes with the displayed
source isomorphism. The primitive chord lift \(U_a\), covering \(-G\), has
target multiplier \(X(p-G)^{-a}\); consequently
\[
 I_\epsilon U_{a,\rm old}I_\epsilon^{-1}
       =\epsilon^{2a}U_{a,\rm new},\qquad
 q_{\rm old}=\epsilon^{-4}q_* .                         \tag{P3K.57.IF.6}
\]
The actual suspension is
\[
 B_{a,b}=2(-1)^{b+1}\lambda q_{\rm old}U_{a,\rm old}^{\,b},
 \qquad ab=2 .
\]
Thus its transported map is exactly
\[
 I_\epsilon B_{a,b}I_\epsilon^{-1}
       =2(-1)^{b+1}\lambda q_*U_{a,\rm new}^{\,b}.       \tag{P3K.57.IF.7}
\]
The exponent \(2ab-4\) is zero. In degree two the Cartier multiplier has the
same scalar factor at source and target; these factors cancel in the
conjugation as well. Therefore \(\lambda_{\rm new}=\lambda_{\rm old}\).



\par\medskip\noindent\textbf{IF.3. The finite root comparison retains its endpoint and lattice period.}\par\nopagebreak[4]

Let the old cap use \(z_{\rm old}^m=T^{-1}\), and the new cap use
\(z_{\rm new}^m=t^{-1}\). On the universal \(\epsilon\)-cover the literal arrow is
\[
 z_{\rm old}=\epsilon^{4/m}z_{\rm new}.                 \tag{P3K.57.IF.8}
\]
The admitted branches are the existing good III-star root
\((a,b,m)=(1,2,4p)\), and reduced III root
\((a,b,m)=(2,1,4p-1)\), with their actual norm sign \(n=\pm1\).
The comparison of minimal-root coordinates is
\[
 \begin{array}{c|cc}
  m&X'_{\rm new}&Y'_{\rm new}\\ \hline
 4p&x'_{\rm old}&y'_{\rm old}\\
 4p-1&\epsilon^{-2/m}x'_{\rm old}&\epsilon^{-3/m}y'_{\rm old}.
 \end{array}                                           \tag{P3K.57.IF.9}
\]
This follows directly by writing \(X'=z_{\rm new}^{2p}x\) and
\(Y'=z_{\rm new}^{3p}y\) and substituting (P3K.57.IF.8). A full positive
\(\epsilon\)-turn has endpoint \(D_{\rm new}^{-4}\). In the reduced case its
factors on \((z_{\rm new},X',Y')\) are
\((\zeta_m^{-4},\zeta_m^{-2},\zeta_m^{-3})\), since \(4p=m+1\).
In the good case the last two factors are one.

The reduced branch retains the actual central nonidentity division section.
On the transported square good-root lattice its chosen logarithm tends to
\((1+i)/2\), and
\[
                   4\log R-\log(4R)=2(1+i).             \tag{P3K.57.IF.10}
\]
The second logarithm is the small identity logarithm. Equation (P3K.57.IF.10)
remains a lattice term in the annular affine cochain. It is not replaced by
zero or multiplied by \(\log\epsilon\) with an assertion of central
holomorphic extension.

The actual full reference overlap is fixed by transporting the prescribed
old ordinary-to-cap overlap at one nonzero parameter, including its integral
delta term and finite line character. The K49 relative source/Hom atlas
comparison then continues that full cocycle on the cap annulus. Its finite
averaging retains the central character and its Cauchy splitting retains the
annular integer. These maps identify the two punctured parameter families.
They do not yet attach a central collision model.

\par\medskip\noindent\textbf{IF.4. Half-loop endpoint at the selected cusp.}\par\nopagebreak[4]

The source equation is even in \(\epsilon\), while \(G\) changes sign. Its
actual endpoint isomorphism is
\[
                  \iota(x,y)=(x,-y).                    \tag{P3K.57.IF.11}
\]
For \(\Theta(QZ)=-Z^{-1}\Theta(Z)\) and
\(\Theta(Z^{-1})=-Z^{-1}\Theta(Z)\), write
\(f_P(Z)=\Theta(Z/P)/\Theta(Z)\). The exact identity
\[
 f_{P^{-a}}(Z^{-1})=P^a f_{P^a}(Z)
\]
gives the selected-cusp line map
\[
                  (Z,V)\longmapsto(Z^{-1},P^a V).       \tag{P3K.57.IF.12}
\]
Here \(P\) is a bounded unit and \(P(-\epsilon)=P^{-1}\). This map squares
to the identity. It sends \(h\) to \(h_{\rm target}^{-1}\), \(g\) to
\(g_{\rm target}\), and preserves the physical delta circle.
Its selected-cusp valuation map is \(\operatorname{diag}(-1,1)\); it sends
the two specified diagonal fans to one another by that actual map.

The cap endpoint of a half \(\epsilon\)-turn is
\(\iota D_{\rm new}^{-2}\). In the reduced case its literal factors from
(P3K.57.IF.9) are \(\zeta_m^{-1}\) on \(X'\) and
\(\exp(-3\pi i/m)\) on \(Y'\); the latter includes the \(y\)-sign of
\(\iota\). Thus the source inversion and cap root rotation are both retained.
The full \(\epsilon\)-turn closes the selected source reflection twice.
In the global rho marking the same arrow is
\(\mathcal R_\rho(Z,V)=(Q/Z,V)\). This is literally
\(h_{\rm target}\circ\mathcal I_\epsilon\), because the target
source period is \((Q,P^{-a})\). Thus the integral height
translation and the factor \(P^a\) in (P3K.57.IF.12) are carried by an
actual period map; neither is dropped from a six-chart closing.

The root coordinate used in the other infinity chart is
\(\zeta=\epsilon/\sqrt{1-\epsilon^2}\), with the square root equal to one
at zero. The ratio is a nonzero unit and has zero winding. A positive
\(\epsilon\)-turn covers a negative double \(v\)-turn. The original
\(u^3=v\) endpoint is further discrete data; it is not closed by
discarding this covering degree.

\par\medskip\noindent\textbf{IF.5. The unselected collision and its actual period orders.}\par\nopagebreak[4]

At \(\epsilon=0\) the cubic discriminant is \(4t^2(1-t)\). The selected
simple zero remains at \(t=1\); the unselected pair meets at \(t=0\).
The exact other two critical values are the analytic branches
\[
 t_\pm=\frac{\epsilon}{8}
 \left\{\epsilon(4-3\epsilon^2)
          \pm\sqrt{(9\epsilon^2-8)^3}\right\}.           \tag{P3K.57.IF.13}
\]
Choose the square root and the labels once. They have distinct first
derivatives \(\pm2i\sqrt2\) and are exchanged by \(\epsilon\mapsto-\epsilon\).
Set \(q_A=t-t_+\), \(q_B=t-t_-\), with the choice agreeing with the marked
sectional small resolution. In its retained theta/Cartier gauge,
\[
 Q=q_Aq_Bu_Q,\qquad P=q_B^{-1}u_P,\qquad
 c_X=q_Bu_X,                                           \tag{P3K.57.IF.14}
\]
where the units are retained as actual functions. The source and
Cartier derivation is the literal chord/theta calculation reproduced
in IF.5a through the unit Weierstrass coordinate map described below.
It retains the degree-two Cartier divisor and its central component.
The leading coefficients in the transported marking are
\[
               u_Q(0)=1/4096,\qquad u_P(0)=64,\qquad u_X(0)=1/8 .
\]

For completeness, a literal source map to the marked local model exists.
Choose either unit root \(\kappa\) of
\[
 \kappa^3-(-2+3\epsilon^2)\kappa-2\epsilon(1-\epsilon^2)
 =(\kappa-\epsilon)(\kappa^2+\epsilon\kappa+2-2\epsilon^2).
\]
Its central value is \(i\sqrt2\) or \(-i\sqrt2\).
Then \(x=\kappa^2 X_0\), \(y=\kappa^3Y_0\),
\(t=\kappa^4T_0\), \(v_0=g/\kappa^3\) sends (P3K.57.IF.1) to
\[
 Y_0^2=X_0^3+(1-2v_0)X_0^2+T_0X_0+v_0^2,\qquad G_0=(0,v_0).
\]
This map also carries the marked small resolution and the canonical
Cartier sections. The whole-zero site is carried to the other, distant
simple discriminant value. It is not changed to one of the colliding sites.
More explicitly, the two source heights are divided by \(\kappa^4\);
the earlier units \(u_Q^0,u_P^0,u_X^0\) become
\(\kappa^{-8}u_Q^0,\kappa^4u_P^0,\kappa^{-2}u_X^0\).
Their previous central values \(1/256,16,-1/4\), together with
\(\kappa^2=-2\), give the displayed constants. This also fixes the signs
of the two node labels by the chosen \(\kappa\) branch.

The actual periods for \(\lambda=\epsilon^N\lambda_0\) are
\[
 h=(Q,P^a),\qquad
 g=(P^{-b},\,2(-1)^{b+1}\lambda_0\epsilon^N q_*c_X^{-2}).
                                                               \tag{P3K.57.IF.15}
\]
Because \(q_*(0,0)=-1\), its collision order is zero. Their divisor matrices
at the two strict transforms, the parameter direction, and their sum are
\[
 P_A=\begin{pmatrix}1&0\\0&0\end{pmatrix},\quad
 P_B=\begin{pmatrix}1&b\\-a&-2\end{pmatrix},\quad
 E=\begin{pmatrix}0&0\\0&N\end{pmatrix},\quad
 H_\infty=\begin{pmatrix}2&b\\-a&N-2\end{pmatrix}.
                                                               \tag{P3K.57.IF.16}
\]
In particular \(\det H_\infty=2N-2\). These are period-order identities,
not a count of fixed strata, links, or characteristic numbers.

\par\smallskip\noindent\textit{IF.5a. Exact source, theta and Cartier derivation of the periods.}\par\nopagebreak[4]

We give the local derivation used in (P3K.57.IF.14)--\allowbreak (P3K.57.IF.15), so no
unpromoted source-zero construction is needed. Work first in the
unit source coordinates \((X_0,Y_0,T_0,v_0)\) just displayed and
write \(a_2=1-2v_0\). Its discriminant and \(c_4\) give
\[
 j^{-1}=\frac{a_2^2T_0^2-4T_0^3-4a_2^3v_0^2-27v_0^4
                         +18a_2T_0v_0^2}
                   {256(a_2^2-3T_0)^3}.                 \tag{P3K.57.IF.16a}
\]
The convergent inverse of the \(j\)-series therefore gives
\(Q=q_1q_2u_Q^0\), with \(u_Q^0(0)=1/256\), on a
smaller two-node polydisc. Here \(q_1,q_2\) are its two
actual discriminant heights, with the chosen small-resolution
label. All the following unit logarithms are fixed on that
full polydisc.

The node map is explicit. Put
\[
 x_c=\frac{-a_2+\sqrt{a_2^2-3T_0}}3,\qquad
 R=(X_0-x_c)\sqrt{a_2+3x_c+X_0-x_c},
\]
using the square roots tending to one. If
\(f(X)=X^3+a_2X^2+T_0X+v_0^2\), then
\(f(X_0)=f(x_c)+R^2\). Thus
\((Y_0-R)(Y_0+R)=f(x_c)\), with
\(f(x_c)=q_1q_2\upsilon\), \(\upsilon(0)=-1/4\).
Dividing one factor by this actual unit gives the literal
node equation \(AB=q_1q_2\). The derivative in \(X_0\)
of \(R\) at the origin is one, so the displayed map is
invertible by the holomorphic inverse theorem.

The convergent Tate sums extend over their node chart
\(Q=Z_0W_0\), with leading terms
\(x_T=Z_0+W_0+\cdots\),
\(y_T=-W_0+\cdots\). Rewriting negative-index terms in
\(W_0\) and \(Q\), and positive-index terms in \(Z_0\)
and \(Q\), proves normal convergence there. After the
specified \(c_4\)-root scaling and Weierstrass translation,
this is the invertible node map above. The actual section
\(G_0=(0,v_0)\) has node factors with linear terms
\(v_0+T_0/2\) and \(v_0-T_0/2\). Its two Tate
coordinates therefore vanish simply on opposite
discriminant branches. Analytic division gives
\[
 Z_0(G_0)=q_1u_G,\quad W_0(G_0)=q_2u'_G,\qquad
 P=Z_0(G_0)/Q=q_2^{-1}u_P^0.                         \tag{P3K.57.IF.16b}
\]
The selected nodal normalization gives the exact value
\(P=(1+\sqrt{1+v_0})/(1-\sqrt{1+v_0})\) on
\(q_1=0\). Since \(q_2=-4v_0+O(v_0^2)\) there,
\(u_P^0(0)=16\). This selects the actual period lift
of the section, including its integral reindexing.

Write \(f_P(Z)=\Theta(Z/P)/\Theta(Z)\). Its divisor
and exact multiplier \(f_P(QZ)=P f_P(Z)\) identify
it with the canonical meromorphic section of the
actual degree-zero line. The source function \(X_0\)
has divisor \([G_0]+[-G_0]-2[O]\), and consequently
\[
 X_0(Z)=c_X^0 f_P(Z)f_{P^{-1}}(Z),\qquad
                     c_X^0=q_2u_X^0.                   \tag{P3K.57.IF.16c}
\]
The coefficient is fixed by evaluating at the ordinary
Tate point \(Z=-1\). The theta shift identity shows
that \(q_2f_P(-1)f_{P^{-1}}(-1)\) is a nonzero
holomorphic unit; so is the actual \(X_0(-1)\).
At the central nodal cubic the retained coordinate
is \(X_0=4Z/(1-Z)^2\), hence \(c_X^0\sim-4/P\)
and \(u_X^0(0)=-1/4\). The unit map by \(\kappa\)
then gives exactly the three constants and functions
in (P3K.57.IF.14).

The primitive source lift covering \(-G\) has target
Cartier multiplier \(X(p-G)^{-1}\). The identity
\(f_P(Z)f_{P^{-1}}(Z/P)=1\) cancels every
\(Z\)-dependent factor in its source and target
Cartier coefficients. On the tensor-power line the
literal resulting map is
\[
          U_a:(Z,V)\longmapsto(P^{-1}Z,c_X^{-a}V).
                                                               \tag{P3K.57.IF.16d}
\]
For \(a=2\), retain the actual Cartier comparison
before using that tensor-power coordinate. With
\(F_2=1-Bz^3-(1-2v_0)z^2\) in the small-resolution
chart of IE.3, the addition formula gives
\[
 r(p)=X_0(p-G_0)=\frac{2v_0F_{21}(p)}{X_0(p)}
                         =\frac{F_2}{z^2}.             \tag{P3K.57.IF.16e}
\]
Its divisor is \([2G_0]+[O]-2[G_0]\). Multiplication
by \(r\) is therefore the actual map
\(\mathcal O([2G_0]-[O])\to
\mathcal O([G_0]-[O])^{\otimes2}\). The last
expression retains its central component and has no
extra factor of \(v_0\). In these unit coordinates the
actual selected height is \(q_0=\kappa^{-4}q_*\).
At \(p'=p-G_0\),
\(F_{21}(p')=r(p')r(p)/(2v_0)\), so its conjugation
of the suspension is the exact identity
\[
 \frac{r(p')}{r(p)}\frac{\lambda q_0}{v_0F_{21}(p')}
                   =\frac{2\lambda q_0}{X_0(p')^2}.
                                                               \tag{P3K.57.IF.16f}
\]
The scalar unit source change is carried on both sides.
This proves (P3K.57.IF.15) for the actual degree-two divisor
line, as well as the degree-one line. It does not
replace that Cartier line by its generic Picard class.

\par\medskip\noindent\textbf{IF.6. The half-loop Cartier arrow through the collision.}\par\nopagebreak[4]

The representative \(P\sim q_B^{-1}\) in (P3K.57.IF.14) requires a different
integral Tate reindexing after the two branches are exchanged. The actual
target representative is
\[
                       P'=(QP)^{-1}.                   \tag{P3K.57.IF.17}
\]
It represents \(-G\), has the required \(q_A^{-1}\) leading order, and
retains the integral period omitted by the bounded selected-cusp expression.
Applying the theta shift identity \(a\) times gives
\[
 f_{(P')^a}(Z^{-1})
   =c_a Z^a f_{P^a}(Z),\qquad
 c_a=(-1)^aQ^{-a(a-1)/2}P^{a(1-a)}.                    \tag{P3K.57.IF.18}
\]
Thus the actual collision line map is
\[
              \mathcal I(Z,V)=(Z^{-1},c_a Z^aV).        \tag{P3K.57.IF.19}
\]
It preserves delta and squares literally to the identity:
\(c_a(P)c_a(P')=1\).
Applying the same two theta identities to the divisor of \(x\) gives
\[
                   c_X'=(QP^2)c_X .                    \tag{P3K.57.IF.20}
\]
Consequently the complete source and suspension periods satisfy
\[
 \mathcal I h=h'^{-1}\mathcal I,\qquad
 \mathcal I g=h'^{-b}g'\mathcal I.                      \tag{P3K.57.IF.21}
\]
For the second equality, its source coordinate has factor
\(P^b=Q^{-b}(P')^{-b}\). Its line coordinate has factor
\(P^{-ab}g_V=(P')^{-ab}(QP^2)^{-2}g_V\), since \(ab=2\).
This proves the full line equality including the scalar, rather than only
the generic Picard-point equality.

Let \(c=a(a-1)/2\). In the ordered valuation lattice
\((Z,V,q_A,q_B)\), (P3K.57.IF.19) and the branch exchange give
\[
 S=\begin{pmatrix}
 -1&0&0&0\\ a&1&-c&c\\0&0&0&1\\0&0&1&0
 \end{pmatrix},\qquad S^2=1,\qquad S\delta=\delta.       \tag{P3K.57.IF.22}
\]
Its fibre block and period relabelling are
\[
 J=\begin{pmatrix}-1&0\\a&1\end{pmatrix},\qquad
 R=\begin{pmatrix}-1&-b\\0&1\end{pmatrix}.
\]
The exact identities
\[
 JP_A=P_BR,\quad JP_B=P_AR,\quad JE=ER,\quad
 JH_\infty=H_\infty R                                  \tag{P3K.57.IF.23}
\]
prove compatibility with every retained period order. The nonconstant
unit factors in \(c_a\) stay in (P3K.57.IF.18); (P3K.57.IF.22) records only their divisors.
Equations (P3K.57.IF.18)--\allowbreak (P3K.57.IF.21) are the actual morphism relating the
selected-cusp and collision half-loop descriptions.

\par\medskip\noindent\textbf{IF.7. The actual origin and the attachment sequence.}\par\nopagebreak[4]

The scalar origin on the ordinary source uses \(\log q_*\). Near the
unselected collision \(q_*\) is a nonzero unit on the full bidisc, so this
logarithm is single valued there. Its \(q_A\) and \(q_B\) continuations do
not insert the extra \(v^{-nk}\) deck found at the selected source-zero
collision. Its parameter endpoint is the extending unit \(q_*^{nk}\).
The finite cap still carries its full norm generator and root arrow
(P3K.57.IF.8)--\allowbreak (P3K.57.IF.10).

The following sections construct the rank-one quotient atlas, its actual finite-cap
attachment, and the real selected-cusp comparison, retaining the physical Euler line
and oriented boundary. The analytic collision filling uses the specified positive
scalar rescaling; IE.6 removes that scalar winding by an actual parameter bordism.

The executed script
\(\texttt{code/part3\_fourfold/source\_infinity\_exact.py}\) checks the
source equations, discriminant branches, selected units, full scalar
cancellation, all root-coordinate exponents, and the complete half-loop
source/line/period identities. Its JSON retains the explicit cap rows
and has global-bordism and return fields false.

\par\medskip\noindent\textbf{IF.8. Literal rank-one sector fans and protected boundary chains.}\par\nopagebreak[4]

All three lower matrices in (P3K.57.IF.16) have rank one. The following is an
actual construction of their cone collections. It is kept in the original
integer lattice, with the physical cocharacter \(\delta=(0,1,0,0)\).
For \(P_0=P_A,P_B,E\), put
\[
 R_0(n)=(P_0n,1,0),\quad R_1(n)=(H_\infty n,0,1),\qquad
 T_j(x,r,s)=(x+(rP_0+sH_\infty)j,r,s).                 \tag{P3K.57.IF.24}
\]
Use the basis \(u=e_1,v=-e_2\) on A, the basis
\(u=e_1,v=e_2-be_1\) on B, and \(u=-e_2,v=e_1\) on E.
In each case \(P_0v=0\). Start with the two labelled
triangles \((0,u,u+v)\), \((0,v,u+v)\). The three
nondegenerate maximal cones in one prism cell are
\[
 \begin{split}
 &\operatorname{cone}(R_0(0),R_1(0),R_1(u),R_1(u+v)),\\
 &\operatorname{cone}(R_0(0),R_1(0),R_1(v),R_1(u+v)),\\
 &\operatorname{cone}(R_0(0),R_0(u),R_1(u),R_1(u+v)).
 \end{split}                                             \tag{P3K.57.IF.24a}
\]
Take all their literal \(T_j\)-translates. The other
three tetrahedra of the two triangular prisms collapse
to the lower-dimensional faces identified below.
In each case \(P_0v=0\). Start with the two triangles
\((0,u,u+v),(0,v,u+v)\) and their three usual staircase tetrahedra
between the \(R_0\) and \(R_1\) copies. Omit their collapsed
four-dimensional cones and retain every ordinary face.

Here is a direct common-face and support proof. Set
\(p=P_0u,\ A_1=H_\infty u,\ B_1=H_\infty v\).
The rational inspection map with basis \((A_1,B_1)\) sends \(p\) to
\((c,d)\), with
\[
 \begin{array}{c|cc}
 &c&d\\ \hline
 A&(N-2)/(2N-2)&-a/(2N-2)\\
 B&(N-2)/(2N-2)&-a/(2N-2)\\
 E&N/(N-1)&bN/(2N-2).
 \end{array}                                           \tag{P3K.57.IF.25}
\]
Every \(c\) is positive for the admitted \(N>2\). At fixed \(r,s>0\)
the two top cones give the square \([0,s]^2\), split along its
diagonal. The third surviving cone gives
\[
          (s,0)+[0,r](c,d)+[0,s](0,1).
\]
Their union is the strip
\[
 0\le x\le L=s+rc,\qquad f(x)\le y\le f(x)+s,\quad
 f(x)=0\ (x\le s),\quad f(x)=(d/c)(x-s)\ (x\ge s).
\]
The actual period vectors are \((L,rd)\) and \((0,s)\).
Reducing \(x\) uniquely modulo \(L\), then reducing
\(y-i rd-f(x-iL)\) uniquely modulo \(s\), assigns every point
off a boundary to exactly one of these strip cells.
The square and parallelogram share their entire vertical edge;
the far edge is the next square's near edge; and the upper
piecewise linear boundary is the lower one translated by \((0,s)\).
Thus all intersections are full common faces. The two collapsed
bottom tetrahedra have the already existing face
\(\operatorname{cone}(R_0(0),R_0(u),R_1(u+v))\), and the collapsed
mixed tetrahedron has
\(\operatorname{cone}(R_0(0),R_1(v),R_1(u+v))\).
No image of a collapsed cone is omitted from the face complex.
This proves the fan and support statements with the exact translations
(P3K.57.IF.24), including the rank-one height boundary.

On the H boundary the A and E constructions give the negative
unit-square triangulation in \(H_\infty^{-1}\) coordinates; B gives
the shear with basis \((e_1,e_2-be_1)\). Their literal polygonal
overlay is a common rational periodic refinement. Intersections
are common faces because their supporting equalities are those of
the two original complexes. It has finitely many cells modulo
\(\mathbb Z^2\). All rays are then replaced by their primitive
integer vectors on the same ray, preserving their true heights.
For \(b=2\) the added H-edge midpoint ray is the primitive multiple
of \((2,N-2,0,2)\).

The A boundary retains every primitive ray \((j,0,1,0)\); the B
boundary retains \((j,-aj,1,0)\). They are exactly the actual source
line chains: on B the literal coordinate \(Y=Z^aV\) has no vanishing
line character along that lower face. The E boundary retains all
primitive rays \((0,j,1,0)\), not only the raw step-\(N\) rays.
Its adjacent coordinates are
\[
       X_j=V\xi_E^{-j},\quad Y_j=\xi_E^{j+1}/V,\qquad
                       X_jY_j=\xi_E .                    \tag{P3K.57.IF.26}
\]
The actual scalar period shifts \(j\) by \(N\), retaining its
nonzero multiplier in the two chart maps.

Extend these boundary refinements by coning already subdivided
faces from their primitive ray sum, face by face in increasing
dimension. This gives common faces and commutes with every
unimodular map (P3K.57.IF.24). The relative regularization then uses a
nonzero lattice point in each singular cone's half-open
parallelepiped, subdividing all occurrences of its minimal face
simultaneously. The incident determinants strictly decrease;
the relative separator prevents two translates of that centre
in one cone. The protected boundary faces are already regular
and hence never contain a resolution centre.

The executable \(\texttt{source\_infinity\_regular\_fan.py}\) is a
separate saved adaptation of the earlier complete algorithm.
Its derivation file records every changed matrix, rank-one
canonical representative, stabilizer and protected-face dimension.
It exports every actual primitive ray, determinant-one cone,
period, face permutation and higher parameter order. Its legacy
tags I,W,E mean the present A,B,E faces; they do not identify A
with the earlier full-rank I face.

\par\smallskip\noindent\textit{IF.8a. Relative separator and terminating integral resolution.}\par\nopagebreak[4]

Here is the complete relative regularization argument for the
present infinity matrices. Protect all the actual A, B and E
lower-face chains listed above. The adjacent primitive rays
have a unit two-dimensional minor; their span lattices are
saturated, with equations respectively \(V=0\),
\(V=-aZ\), and \(Z=0\). Hence every protected cone is
regular in its actual lattice. In each sector the protected
subfan is exactly its lower face \(s=0\). A nonprotected
ray has positive H height. Its period stabilizer is zero:
\[
 \begin{aligned}
 \det(rP_A+sH_\infty)&=\det(rP_B+sH_\infty)\\
                    &=(N-2)rs+(2N-2)s^2,\\
 \det(rE+sH_\infty)&=2Nrs+(2N-2)s^2.
 \end{aligned}                                         \tag{P3K.57.IF.26a}
\]
Both are positive for \(r\ge0,s>0,N>2\). Thus
\(T_\nu(x,r,s)=(x,r,s)\) on such a ray implies
\(\nu=0\). The lower-face kernel stabilizers fix those
faces pointwise; their actual analytic unit actions are
retained in IF.9 and are not replaced by identity maps.

First take the relative barycentric subdivision outside
the protected subfan. Process old faces in increasing
dimension, insert the primitive ray on the sum of their
primitive generators, and cone the already subdivided
boundary from that ray. A convex transverse section gives
the literal radial coverage and common-face map. The
same choices are made on the common H faces in all three
sectors. Every period map is unimodular, preserves face
dimension, and takes a generator sum to the corresponding
generator sum; hence this subdivision is exactly equivariant.

Each resulting cone consists of a protected face and a
strictly increasing chain of nonprotected old faces. It
therefore contains at most one ray in any nonprotected
period orbit: two rays in an orbit would arise from old
faces of the same dimension, whereas a strict face chain
contains only one of that dimension. An inserted ray is
interior to its named face, so equality of rays also gives
equality of their named faces. Every singular face contains
a nonprotected ray, because a cone generated entirely by
protected rays is a face of one of the regular protected
cones. Two distinct period translates of a singular face
cannot lie in one cone: their nonprotected rays would
violate this orbit separation, and the ray stabilizer is
trivial by (P3K.57.IF.26a).

For a singular maximal cone let \(B\) be its primitive
column matrix and \(D=|\det B|>1\). Some column
\(B^{-1}e_j\) is nonintegral, since otherwise \(B^{-1}\)
would be integral and its determinant would force \(D=1\).
Take its coordinatewise fractional part \(\alpha\) and set
\[
 p=B\alpha=e_j-B\lfloor B^{-1}e_j\rfloor,
 \qquad 0\le\alpha_i<1.                              \tag{P3K.57.IF.26b}
\]
This is a nonzero integral point. Divide \(p\) and every
\(\alpha_i\) by the positive gcd of its coordinates, so
\(p\) is primitive and the same strict inequalities hold.
Its minimal face \(F\), generated by the rays with
\(\alpha_i>0\), is singular. Indeed on a regular face
an integral vector with all these coefficients in
\([0,1)\) is zero in the saturated span lattice, which
contradicts \(p\ne0\). Thus \(F\) is outside the
protected subfan and has trivial period stabilizer.

Insert \(T_\nu p\) simultaneously on every translated
face \(T_\nu F\). Orbit separation proves that no old
cone receives two distinct centres from this orbit.
For a point \(x=\sum t_i r_i\) of a containing cone,
subtract \(\tau p\), where
\(\tau=\min_{\alpha_i>0}t_i/\alpha_i\). The remaining
coefficients are nonnegative and at least one is zero,
which gives precisely a cone of the star subdivision.
Ties are precisely the shared faces; the same formula
on an old shared face proves compatibility. For each
affected maximal cone and each replaced ray,
\[
                        D_{
m new}=\alpha_i D_{
m old},
                 \qquad 0<\alpha_i<1.                  \tag{P3K.57.IF.26c}
\]
These are positive integers. Every affected determinant
strictly decreases, and no protected cone is subdivided.

Orbit separation survives the star: a new cone lies in
an old cone and contains at most one new ray of the chosen
orbit. This new orbit was not an old ray orbit, since its
point was interior to a face of dimension at least two.
Its positive H height gives trivial stabilizer again.
There are initially finitely many cone orbits, and every
star introduces finitely many new ones. At each stage
choose a maximal-determinant cone orbit. The pair
\((D_{\max},\#\{\text{maximal orbits of determinant }D_{\max}\})\)
strictly decreases lexicographically: no new orbit has
the old maximal determinant. There are finitely many
stages at each first entry and finitely many decreases
of that entry. Hence the process terminates with every
maximal determinant one. If the centre lies on H, perform
the identical star in all incident sectors; formula
(P3K.57.IF.26c) applies to each. This proves a compatible
period-equivariant regular refinement relative to all
three actual chains, in the original integral lattice,
with every higher parameter height retained.

\par\medskip\noindent\textbf{IF.9. Exact face kernels and the analytic quotient.}\par\nopagebreak[4]

Put \(C=2(-1)^{b+1}\lambda_0\). On the A face the kernel period is
\(g\). Its transverse line coordinate is \(V\), and its multiplier is
\[
                  C\epsilon^N q_*q_B^{-2}u_X^{-2}.       \tag{P3K.57.IF.27}
\]
On the B face the kernel is \(k=g h^{-b}\). Its source multiplier
and its multiplier on the literal transverse coordinate \(Y=Z^aV\)
are
\[
 (q_Au_Qu_P)^{-b},\qquad
          C\epsilon^N q_*q_A^{-2}u_Q^{-2}u_P^{-4}u_X^{-2}.       \tag{P3K.57.IF.28}
\]
These equalities follow directly from (P3K.57.IF.15): the second line is
\(g_V/(QP^2)^2\). Since the two critical values differ by
\(\epsilon\) times a nonzero unit, both raw transverse kernel
multipliers have order \(N-2\) at their corresponding strict face.
All constants and source multipliers remain in (P3K.57.IF.27)--\allowbreak (P3K.57.IF.28).
They are not replaced by a formal valuation.

The raw multiplier order \(N-2\) must be combined with the actual nodal-line metric and source reindexing. The complete operator bounds below retain losses three and four, respectively, and the geometric proof uses the safe orders \(N=8\) and \(N=6\).

At E the source period \(h\) retains the actual compact elliptic
source fibre and its line holonomy. The scalar period has order \(N\)
and acts on the primitive chain (P3K.57.IF.26). Near H, both previously
kernel directions have nonzero image under the full-rank matrix
\(H_\infty\). Thus the period separation in the full-rank cones
controls the closures of both face norm shells; no additional
kernel is lost at that overlap.

Here is the complete analytic quotient argument, including the
norm bound for the actual nodal lines. The analytic charts are
the affine semigroup charts of the regular fan, pulled back by
the stated base-blowup maps. The full period indexed by
\(\nu\in\mathbb Z^2\) has lattice map \(T_\nu\) from
(P3K.57.IF.24), followed by the actual extending base units from
(P3K.57.IF.15). Multiplication by these units acts on every character
by its literal integer powers, and therefore extends to the
corresponding semigroup chart. The units form a homomorphism
in \(\nu\), since they are products of the two commuting
periods. This proves the group law, its inverse, and the
common-H overlap law on the analytic atlas.

For a nonzero base value, let \(r,s\) be the positive logarithmic
radii of the two sector coordinates. The full real matrix of
the logarithmic magnitudes of its two actual periods is
\[
                       L=-rP_0-sH_\infty+C,             \tag{P3K.57.IF.28a}
\]
where \(C\) is bounded; the fixed \(\lambda_0\) belongs
to that bound. By (P3K.57.IF.26a) the determinant of the leading
matrix is bounded below by a positive multiple of
\(s(r+s)\), and its norm is bounded above by a multiple
of \(r+s\). Its least singular value is therefore at
least a positive multiple of \(s\). Shrinking the corner
polydisc makes the bounded perturbation harmless. Use the
actual parallelogram \(L([0,1]^2)\), with its compact
angular torus. It differs from the leading parallelogram
by a uniformly bounded set.

The strip construction in IF.8 places that set in finitely
many adjacent chart types, uniformly in \(r/s\). Indeed
its width is \(s+rc\), its height is \(s\), and its
periods are \((s+rc,rd)\), \((0,s)\). Compare the
linear function \(rd\,x/(s+rc)\) with the two pieces
of \(f(x)\). Their difference divided by \(s\) is
bounded uniformly, so only a bounded number of vertical
translates is required. A bounded perturbation changes
both assigned strip integers by a bounded amount.
Consequently one fixed finite collection of adjacent
cone charts covers the quotient. Its monomial coordinate
magnitudes are uniformly bounded above and its angles
lie in compact tori. Closing those regions, including
their toric faces, gives compact chart sets over a
closed smaller base polydisc. Away from the three
marked H points the same argument uses the full-rank
matrix \(H_\infty\) alone.

At the strict B face, restrict to
\(|q_B|\le\eta|q_A|\), with a fixed small \(\eta\).
Then \(|\epsilon|\) is comparable with \(|q_A|\).
The nodal line in the character \(Y=Z^aV\) has source
period multiplier \(c=(QP)^a=(q_Au_Qu_P)^a\).
On \(|Q|\le|Z|\le1\), interpolate a positive metric
between its two actual endpoint values. Its ratio is
bounded by \(C|q_A|^{-a}\). The same interpolation on
the plumbing annuli extends to the nodal normalization
with the prescribed gluing at both node preimages.
The source multiplier of the kernel is bounded by
\(C|q_A|^{-b}\). Reindexing its image to the Tate
fundamental annulus uses at most two source periods
for \((a,b)=(1,2)\), and at most one for \((2,1)\).
Each period changes the line norm by at most
\(C|q_A|^{-a}\). Combining these bounds with the
raw multiplier of order \(N-2\) gives
\[
 \|k\|\le
 \begin{cases}
 C|\lambda_0|\,|\epsilon|^{N-5},&(a,b)=(1,2),\\
 C|\lambda_0|\,|\epsilon|^{N-6},&(a,b)=(2,1).
 \end{cases}                                            \tag{P3K.57.IF.28b}
\]
At A restrict instead to \(|q_A|\le\eta|q_B|\).
Here \(|\epsilon|\) is comparable with \(|q_B|\),
the literal line character is \(V\), and its source
period multiplier is \(P^a\), of norm comparable
with \(|q_B|^{-a}\). Choose its endpoint metric
with the prescribed ratio \(|P|^a\), and interpolate
on the source annulus and plumbing charts. Its
maximum-to-minimum norm ratio is at most
\(C|q_B|^{-a}\), including both preimages of the
nodal point. The source part of \(g\) is
\(P^{-b}\), of norm comparable with \(|q_B|^b\).
Since \(|Q|\le C|q_B|^2\), returning the image of
\(|Q|\le|Z|\le1\) to that same annulus uses at most
two source periods for \(b=2\) and at most one for
\(b=1\), after fixing \(\eta\) and shrinking.
Each positive or negative reindex changes the line
norm by at most \(C|q_B|^{-a}\). The raw multiplier
has order \(N-2\); adding this metric loss and those
reindex losses gives exactly (P3K.57.IF.28b) for \(g\).

These estimates include all their face overlaps.
For fixed \(\eta<1\), the strict A and B regions
are disjoint; the intervening sector is covered by
the full-rank H parallelograms already constructed.
On each strict face its middle annulus and two node
charts use the same endpoint metric and transition
multiplier, so no further norm ratio is introduced
on their overlaps. Away from H the relevant
source/line maps lie in compact families and have
a common finite norm bound. At E the kernel is
the actual source Tate period, with its complete
elliptic fundamental annulus and its line holonomy.
The common H charts retain the nonzero images of
all these kernel vectors, which controls the
closures of their norm shells as described below.

We use the analytic construction at the two executed
orders \(N=8\) and \(N=6\), respectively. In the first
case (P3K.57.IF.28b) tends to zero; in the second it is
uniformly less than one after choosing \(|\lambda_0|\)
small. For an arbitrary admitted cap, take a sufficiently
large multiple of its order; IE.6 compares these orders
by an actual scalar pair of pants. The additional
\(N=3,4\) fan executions are finite arithmetic checks,
not an assertion of this norm bound at those orders.
The actual line-norm shells prove freeness and proper
discontinuity of the retained kernels on the nonzero
nodal lines. On E use the actual Tate shell for
\(Z\mapsto QZ\) together with the primitive chain.
Near H the two kernel vectors have nonzero H images,
so an unbounded sequence of their translates cannot
accumulate in a fixed H-cone chart. The full-rank
separation controls the closures of the face shells.

Finally a nonzero period cannot stabilize a bounded
cell on a full-rank face, and only finitely many
translates of a compact chart meet it. On a rank-one
face the only possible stabilizer is the specified
kernel, whose strict norm contraction has the same
finite-intersection property on a compact shell.
The equivalence relation is locally a finite union
of closed graphs on these compact covers. Its quotient
is Hausdorff, and inverse images of compact smaller
base sets are compact. On A and B its boundary charts
are precisely the actual Tate source chain and
Cartier line with their retained kernel maps; on E
they are (P3K.57.IF.26). Thus the punctured-parameter part
is the prescribed family under the displayed theta
and Cartier maps, with the stated boundary charts.

The global E-chain monodromy still carries its actual source-base
data. A small loop around the unselected collision has zero winding
of \(q_*\), so it fixes all \(N\) chain labels. A loop around the
separate selected value shifts labels by \(-nk\), and the inverse
cap cut at infinity shifts them by \(+nk\), for \(N=km\).
The latter two have \(k\) cycles of length \(m\). This is a
compatible list of branch locations; it does not itself construct
the missing selected-cusp attachment.

The real joining of this partial filling to the selected scalar trace is
constructed in IE.1--6, with its physical Euler line and normal orientation.

\par\medskip\noindent\textbf{IF.10. The changed-origin partial attachment away from the selected disc.}\par\nopagebreak[4]

Fix disjoint source-base discs around the unselected collision at \(t=0\),
the selected value \(t=t_*(\epsilon)\), and the cap at infinity. The first
and last closures are disjoint from the selected disc for all sufficiently
small \(\epsilon\). Put \(N=km\), retain the actual \(n=\pm1\), and use
\[
 \tau=\frac{\log\epsilon}{2\pi i},\quad
 s_*=\frac{\log q_*}{2\pi i},\quad
 s_\infty=\frac{\log q_\infty}{2\pi i},\quad
 h=1-t_*(\epsilon)q_\infty,\qquad
 s_*=-s_\infty+\frac{\log h}{2\pi i}.                   \tag{P3K.57.IF.29}
\]
On the unselected collision bidisc \(q_*\) is a nonzero unit, so \(s_*\)
has one holomorphic branch on that full bidisc. On the cap polydisc \(h\)
is likewise a nonzero unit, with its chosen reference logarithm.

On the universal parameter cover the actual ordinary-to-cap map and
the actual positive parameter-deck square are
\[
 F_\tau=F_0+nk\,\tau s_\infty\delta,\qquad
 R_{\rm ord}=q_*^{nk},\quad R_{\rm cap}=h^{nk},\qquad
 F_{\tau+1}R_{\rm ord}=R_{\rm cap}F_\tau.               \tag{P3K.57.IF.30}
\]
These are the admitted cap's norm and cochain formulas, with the current
source and the transported full \(F_0\) of IF.3. The source and line
characters are not changed. Define
\[
 o_{\rm ord}=nk\,\tau s_*\delta,\qquad
 o_{\rm cap}=nk\,\tau\frac{\log h}{2\pi i}\delta,\qquad
 \eta=\zeta-o_{\rm ord}.
\]
The exact overlap is
\[
                F_\tau+o_{\rm ord}-o_{\rm cap}=F_0.    \tag{P3K.57.IF.31}
\]
After a positive parameter turn the old ordinary chart is identified
by \(R_{\rm ord}\). Its change is exactly the change of \(o_{\rm ord}\);
therefore \(\eta\) has identity parameter deck. The cap coordinates
have the same property. This retains the closing arrow: neither
\(R_{\rm ord}\) nor \(R_{\rm cap}\) is silently set equal to one in
the old chart.

Every source period and line lift commutes with these base-only scalar
translations. Hence in the new ordinary coordinate the full periods
are exactly (P3K.57.IF.15), including all their units. Moreover positive
\(q_A\) or \(q_B\) continuation in the collision bidisc leaves \(s_*\)
unchanged. Thus the new collision coordinate has no additional
source-base chain deck there. The fan of IF.8 may therefore be
constructed directly in \(\eta\), and its quotient of IF.9 has the
pure parameter deck just proved. On its punctured parameter part
this is an actual isomorphism to the prescribed family, using
(P3K.57.IF.30), not an assertion that a fractional scalar power extends in
the old \(V\) coordinate.

The cap chain is constructed in \(\zeta_{\rm cap}-o_{\rm cap}\).
The actual primitive generator is
\[
 H=A^uB^w=\epsilon^k r_0H_0,\qquad nu+mw=1,
 \qquad H^m=B,\quad H^n=A ,
\]
with the actual admitted source and line maps. The displayed Bezout
equation is the literal exponent identity for \(A=H^n,B=H^m\).
Thus its node chart shift is
\(k\), and \(A\) has signed shift \(nk\). The source and finite
root maps in \(H_0\) are retained. The chain equations and all line
frame transitions are (P3K.57.IF.26), with \(H\) acting by that shift and
by its actual source/line map. A nonzero iterate changes the index,
so the action is free and properly discontinuous also at
\(\epsilon=0\).

On the cap annulus positive \(q_\infty\) continuation changes
\(s_*\) by \(-1\); the new ordinary vertical coordinate consequently
has the factor \(\epsilon^{nk}\). This is exactly \(A\)'s chain
shift and its retained \(A_0\) arrow. The cap unit \(h^{nk}\) acts
on a node chart by
\((X_j,Y_j)\mapsto(h^{nk}X_j,h^{-nk}Y_j)\); it is an extending
unit even though its central value is now \(1-q_\infty\), rather
than the constant one of the source-zero calculation. Its origin
change has been retained in (P3K.57.IF.30)--\allowbreak (P3K.57.IF.31).

On each remaining ordinary source-base chart, away from the selected
disc and these two special pieces, the source and Cartier line are
smooth over the full parameter disc. Use their full chain with the
actual source/line period. On angular overlaps the affine cochains
are exactly those already prescribed by \(F_0\); the only nonzero
integral base continuation is the signed cap shift just displayed.
The collision overlap has zero such continuation. Consequently
the ordinary chains, the collision fan in \(\eta\), and the cap
chain glue by literal extending source/line maps. This constructs
the partial family over the complement of the selected source-base
disc.

The relative regularization is by primitive determinant-one cones,
and the protected face charts are the actual smooth source/chain
charts. The total complex fourfold is therefore smooth. The map
to the epsilon disc may have a singular central fibre; smoothness
here refers to the total space and does not replace that map by a
submersion.

Properness over a closed smaller parameter disc follows from the
finite compact fundamental covers already described: collision
cones and face shells in IF.9, finitely many node/middle chain
charts on every ordinary compact source patch, and one finite
fundamental chain segment for \(H\) on the cap. A finite source-base
cover suffices. Every overlap is the actual biholomorphic map
(P3K.57.IF.31), so these compact covers produce a compact inverse image.
The only retained real boundary is over the boundary circle of the
deleted selected disc.

This partial attachment supplies the exact source-base boundary used in
IE.1--6. Those sections construct its real comparison with the selected
scalar trace, with the physical Euler line and its normal orientation.

\par\medskip\noindent\textbf{IF.11. Executed local fixed-stratum data.}\par\nopagebreak[4]

The two independently saved regular fans have respectively
\[
 \begin{array}{c|r|r|rrr}
 (a,b,N)&\text{checks}&\text{stars}&A&B&E\\ \hline
 (1,2,8)&220943&1289&1838&1838&1818\\
 (2,1,6)&222609&1381&1748&1768&2021
 \end{array}
\]
maximal cone orbits. Every exported ordered cone has determinant \(+1\).
The fixed components are determined by the minimal face containing
the actual \(\delta\), with every recorded period and vertex
permutation used. In particular neither rank-one A nor rank-one B
contains a fixed component on its protected boundary.

The compact interior and H-normal calculations use the local
normal-link convention and formula (P3K.57.IE.16). On H the full unit
factors, with \(w=q_B/q_A\), are
\[
 h:(Z,V)\mapsto(wZ,w^{-a}V),\qquad
 g:(Z,V)\mapsto(w^bZ,w^{-2}(1-w)^NV)
\]
after the \(H_\infty\) normal powers are separated; the nonzero
constants remain part of the atlas. Thus the three normal-degree
orders at \(w=0,1,\infty\) are the same literal expressions
(P3K.57.IE.22), with the actual \(H_\infty\) used when unwrapping
the recorded periods.

The executed values, verified component by component using two
independent auxiliary circles, are
\[
 \begin{array}{c|rrr}
 (a,b,N)&\text{compact interior}&H&\text{sum}\\ \hline
 (1,2,8)&9333/16&-9301/16&2\\
 (2,1,6)&739&-734&5 .
 \end{array}                                           \tag{P3K.57.IF.32}
\]
All H surface tangent Euler numbers, canonical squares,
signatures and holomorphic Euler numbers satisfy their exact
integrality checks.

Each E fixed surface has the actual unmodified local I2 cycle
with intersection matrix \(\left(\begin{smallmatrix}-2&2\\2&-2\end{smallmatrix}\right)\).
The relative determinant marking \(Z^aV^2/\xi_E^{2j+1}\)
has divisor zero on both components, as the unwrapped cone
characters show. The local tangent marking
\(\Omega=d\log Z\wedge dt\) gives relative
\(p_1=-4\). Indeed the actual source-period change adds
\(d\log Q(t)\wedge dt=0\), so \(\Omega\) descends.
On a node chart \(xy=t\), it is \(dx\wedge dy\).
The vertical vector is \(x\partial_x-y\partial_y\);
it has one simple complex zero at each of the two nodes.
Together with the boundary vector normalized by
\(\Omega(W,H)=1\), this gives relative \(c_1=0\)
and \(c_2=2\), hence \(p_1=c_1^2-2c_2=-4\).
Thus the annular normal-link residue is \(-1/6\) per
E surface. These local numbers are retained with their
boundary frames. IE.3 constructs the actual E double,
and IE.4--6 assemble its completed-family contribution.

All programs and their full exact outputs are saved with the prefix
\path{code/part3_fourfold/source_infinity_}.
The source-family return uses the complete IE.1--7 comparison,
including its actual boundary frames and parameter orientation.

\par\medskip\noindent\textbf{IF.12. Full old-v / new-epsilon annular family comparison.}\par\nopagebreak[4]

Fix an integer root-scalar winding \(j\), and use
\[
 r(\epsilon)=\epsilon^j u(\epsilon)r_0,\qquad
 \lambda=r^m ,
\]
where the specified nonzero unit \(u\) has an extending
logarithm. The comparison in this section is over the
punctured epsilon annulus carrying this actual path.
It does not assume a trivial Isom torsor over a product
of two independent punctured parameter discs.

The source map (P3K.57.IF.1), the canonical divisor-line map,
and (P3K.57.IF.6)--\allowbreak (P3K.57.IF.7) identify the full ordinary suspension.
The selected Tate \(Q,P\) agree under this source
isomorphism. Their physical heights differ by
\(q_{\rm old}=\epsilon^{-4}q_*\), an actual single-valued
integer-power unit on this annulus. Its powers act on
the toric character charts and preserve the chosen
periodic fan and physical delta. Its winding is retained.

On each simply connected parameter chart use the actual
root map (P3K.57.IF.8)--\allowbreak (P3K.57.IF.9). It carries the marked source
section, the resolution ideals and the corrected Cartier
line. The finite source/line map has the form \(A=LD\),
where \(L\) covers \(-nR\). Its central commutator
character is one: in the good branch \(R=O\) centrally
and \(L\) is a scalar; in the reduced branch the
corrected root line is trivial on the compact reduced
III fibre, so the multipliers of \(D,L\) are constants,
and their source maps commute. The finite averaging
conjugation of the actual norm-one unit gives
\[
 DLD^{-1}=L,\qquad A^4D^{-4}=L^4 .
\]
The source part of \(L^4\) specializes to the identity
in both branches. Divide its central scalar and take
its small extending Lie logarithm. The actual annular
difference from four times the logarithm of \(L\)
retains the period (P3K.57.IF.10), with its full line factor
and signed \(n\). No multiple of that period is
asserted to be an extending reduced-root logarithm.

Choose one reference \(\epsilon_*\) and transport the
entire old ordinary-to-cap overlap through these
source/root/Cartier arrows, including \(d\delta\),
the chosen root cut, norm and finite character.
Use this full map as the reference of the new
\(F_0\). On a fixed new cap annulus the source and
root families are smooth through epsilon zero.
Restriction to a chosen root sheet identifies
the ordinary line quotient with the cap quotient:
an element \(A^hB^\ell\) fixing that sheet has
\(h=mk\), so is exactly \(B^{nk+\ell}\).
Thus the local comparisons are literal source
and line maps, and their sheet changes are the
retained powers of \(A\).

Correct those maps at \(\epsilon_*\) to the full
reference overlap. On two angular base charts
their overlap translations have Lie logarithms
zero at that reference. Lift along the simply
connected epsilon disc. Every ordinary and
finite-group discrepancy is an integral period
and is zero at the reference, hence stays zero.
The fixed-contour Cauchy splitting in the base
variable therefore glues the actual local maps
to \(F_0\), holomorphically in epsilon. This
retains its full annular integer and the
doubled reduced period. It is the explicit
Hom/source-period torsor construction of
P3K.49, RCZ.5, with the present displayed root maps
as its input.

Now compare the two completed families on a
small simply connected parameter chart.
Their total ordinary/root/cusp lattice
two-cocycles agree under the actual source
period identification: they agree at the
reference, and their integral discrepancies
are continuous in the parameter. The reduced
component translations are first removed by
their actual suspension powers. The remaining
Lie bundle has the literal coarse-base exact
sequence
\[
     0\longrightarrow\mathcal O\delta
       \longrightarrow V\longrightarrow\mathcal O(-1)
       \longrightarrow0 .
\]
The source dualizing vector and physical delta
give its relative frames. Split the extension
and its one-cocycles by the fixed-circle
Cauchy/Laurent operators for \(\mathcal O(1)\),
\(\mathcal O\), and \(\mathcal O(-1)\).
Finite averaging preserves the full root
cocycles. Exponentiation then gives actual
local family isomorphisms extending on the
chosen cusp and root charts.

On two parameter charts their ratio is a
global extending translation. The actual
source/rational-divisor endpoint argument
of P3K.49, RCZ.4 identifies that group with the
physical \(\mathcal O^*\), including all
finite-kernel choices. To glue the local
maps, take logarithms \(\ell\) of their
scalar ratio on the two angular overlap
components, choose \(\chi_++\chi_-=1\),
and set \(b_+=\chi_-\ell,b_-=-\chi_+\ell\).
Their common \(\bar\partial\) is solved on
a smaller annulus by
\[
 v(\epsilon)=\frac1\pi\int_{\mathbb C}
       \frac{\widetilde f(w)}{\epsilon-w}\,dA(w).
\]
Then \(h_\pm=\exp(b_\pm-v)\) have ratio
equal to the actual scalar transition,
and \(h_\pm^{-1}I_\pm\) glue literally.
The cutoff is supported in the larger
annulus and equals one on the loop.
This construction retains both overlap
components and any integral scalar winding.

There is a concrete forward-endpoint
factor before this last scalar splitting.
For the root winding \(j\), the old and
new ordinary forward decks have ratio
\[
      q_{\rm old}^{nj}/q_*^{nj}
                         =\epsilon^{-4nj}.
\]
With \(\tau=\log\epsilon/(2\pi i)\), retain
the actual compensating cochain
\[
 c_0(\tau)=\exp\{2\pi i\,2nj\,\tau(\tau-1)\},
 \qquad
 c_0(\tau+1)/c_0(\tau)=\epsilon^{4nj}.                 \tag{P3K.57.IF.33}
\]
This is not declared single valued by
itself. It belongs to the forward
parameter-deck square, alongside the
\(D^{-4}\) root endpoint and the
retained period above.

All constructed maps and inverses are
source/Cartier coordinate maps,
extending translations on their stated
charts, or physical scalars. Thus they
give a holomorphic delta-equivariant
family isomorphism on the epsilon
annulus, carrying the actual line,
fan, normal circle and closing arrows.
A connection transported through
this family map identifies the physical
returns by conjugacy; there is no
loop reversal or numerical class division.

For the old global scalar
\(r=\epsilon^{-2k}(1-\epsilon^2)^kr_0\),
one has \(j=-2k\). Formula (P3K.57.IF.33) is
\(\exp\{-2\pi i\,4nk\,\tau(\tau-1)\}\),
whose one-step ratio is
\(\epsilon^{-8nk}\). The nonzero unit
\((1-\epsilon^2)^k\) is retained and
can be homotoped to its reference value
through its extending logarithm.
The local rescaling
\(r=\epsilon^{k_\infty}r_\infty\)
is instead \(j=k_\infty\). Their
root-winding difference is
\(-2k-k_\infty\), as in RG.5.
These are two explicitly compared
loops, not two notations for one scalar.


\par\medskip\noindent\textbf{IE.1. Two actual collision fillings and their boundary families.}\par\nopagebreak[4]

Choose a small closed source-base disc \(U\) around \(t=0\), disjoint
from the selected value and the cap. Its boundary annulus contains no
singular source fibre for any small \(\epsilon\). Fix a nonzero
reference \(\epsilon_*\). Let \(S_\epsilon\) be the actual minimal
source over \(U\), with its marked small resolution at \(\epsilon=0\),
and retain \(M_a=\mathcal O(aG-O)\), \(a=1,2\).

The first filling \(W_{\rm c}(N)\) is the regularized collision family
IF.8--10, with \(\lambda=\epsilon^N\lambda_0\). After its
determinant-one regularization the total fourfold is smooth; the
function \(\epsilon\) may have a singular central fibre. Its actual
source/line periods are (P3K.57.IF.15), its physical circle is \(\delta\),
and its protected boundary is the primitive source-line/parameter
chain.

The reference filling \(W_{\rm r}(N)\) uses the fixed source
\(S_{\epsilon_*}\), its fixed actual line \(M_a\), and the same scalar
\(\lambda=\epsilon^N\lambda_0\). No selected zero lies in \(U\).
Its actual suspension lift \(B_0\) is therefore invertible over
\(U\), including both ordinary I1 fibres. Use the full chain
\[
 X_j=V\epsilon^{-j},\qquad
 Y_j=\epsilon^{j+1}/V,\qquad X_jY_j=\epsilon ,
 \quad j\in\mathbb Z.                                  \tag{P3K.57.IE.1}
\]
The suspension carries chart \(j\) to \(j+N\), with \(B_0\) on
its source and line frames. A nonzero iterate changes the chain
index; a segment of \(N\) node/middle charts is a compact
fundamental cover over the compact source patch. Thus
\(W_{\rm r}(N)\) is a smooth fourfold with exactly \(N\) fixed
source surfaces in its central chain, of normal weights \(+1,-1\).

On the boundary annulus the varying source, Cartier line,
translation and line lift all extend over the full
\(\epsilon\)-disc. They are the full source/line maps, not just
their limiting Picard points. The changed-origin coordinate
of IF.10 has pure parameter deck, and its source-base
continuation on this annulus is zero because \(q_*\) is a unit
on \(U\).

Here is the actual smooth boundary transport. Introduce
\(s\in[0,1]\), use source parameter
\(\epsilon_s=(1-s)\epsilon+s\epsilon_*\), and keep the
degeneration coordinate \(\epsilon\) and scalar
\(\epsilon^N\lambda_0\) fixed. The boundary source fibres
stay smooth. Continue their actual line and suspension
using the source/Cartier expressions. On each chain chart
the equation is always \(XY=\epsilon\). There are local
lifts of \(\partial_s\) keeping \(\epsilon\) and the
source-base boundary coordinate fixed: take the
source-family chart lift and its actual line frame,
with zero \(X,Y\) derivatives. Under a unit frame change,
carry its derivative into the two opposite normal
coordinates. These give lifts on the actual quotient
charts.

Use a finite partition of unity on the compact quotient
collar to combine the lifts, then average over the physical
circle. The resulting field still has \(ds=1\),
\(d\epsilon=0\), and fixed source-base boundary coordinate.
Its compact fundamental cover makes the flow complete for
\(0\le s\le1\). Its time-one map is an actual equivariant
boundary diffeomorphism
\[
        \Psi:\partial_U W_{\rm c}(N)
                   \longrightarrow\partial_U W_{\rm r}(N).   \tag{P3K.57.IE.2}
\]
It retains the full source/line units and \(j\)-label
transport. The flow is in the chain atlas, not a
translation logarithm across the interior collision.

Glue the two fillings along (P3K.57.IE.2), reversing the reference
orientation and rounding the common corners. The result
\(W_{\rm rel}(N)\) is a compact oriented smooth
eight-dimensional delta-manifold. Its outer boundary is
the closed relative mapping-torus piece. Its fixed set
consists of the compact interior/H components of
\(W_{\rm c}\), and \(N\) closed doubles of its E patches
with the oppositely oriented reference patches. Their
normal bundles are glued by the actual derivative of
\(\Psi\).

\par\medskip\noindent\textbf{IE.2. Exterior transport and selected-cusp cancellation.}\par\nopagebreak[4]

The same source deformation extends on the complement
of \(U\), including the actual cap root. Its division
section, finite character, norm lift and full annular
integral cochain are those of IF.3 and IF.10. On the
selected I1 collar the actual normalization is (P3K.57.IF.2);
the Tate, Picard and Hom units (P3K.57.IF.3)--\allowbreak (P3K.57.IF.5) have
extending logarithms and zero source-parameter winding.
Its Cox products continue through \(\epsilon_s\), with
the same scalar parameter and physical circle.

On the outer circle the completed threefolds are smooth
and proper. Lift \(\partial_s\) as above, prescribing
(P3K.57.IE.2) on a collar of \(\partial U\). Extend the lift
using a partition of unity and circle averaging.
Near the selected fixed sphere prescribe the literal
Cox-coordinate unit map. Its normal derivative extends
through \(s\) and the source-parameter disc, so it
transports the actual Hopf Euler line and normal
framing. The resulting exterior diffeomorphism agrees
with (P3K.57.IE.2).

The two selected fixed-sphere traces and their normal
products are thereby identified, including the positive
physical orbit. Remove the product tubes and perform
the same conifold replacement on both sides. The
identification extends over \(D^3\times S^2\), with
its transported meridian and Hopf Euler line. The
remaining exterior is a product bordism. Therefore
the difference of the two completed Jin six-manifolds
is the relative collision boundary of \(W_{\rm rel}\),
with the parameter-convention sign in IE.5.
This cancels actual selected scalar traces; it does
not supply a holomorphic full selected corner.

\par\medskip\noindent\textbf{IE.3. The E doubles and actual Cartier class.}\par\nopagebreak[4]

The marked source small resolution in the unit
Weierstrass coordinates of IF.5 is
\[
 A=Dz,\quad C=Bz,\quad X=Bz,\quad
 Y=(B+Dz)/2,\quad v_0=(B-Dz)/2 .
\]
Here \(A=Y-v_0\), \(B=Y+v_0\), \(C=X\), and
\(D=T_0+X^2+(1-2v_0)X\). The complementary chart is
\[
 D=Aw,\quad B=Cw,\quad w=z^{-1},\quad
 v_0=(Cw-A)/2.
\]
Thus \(dv_0/dB=1/2\) in the first chart and
\(dv_0/dA=-1/2\) in the second. These are actual
submersion charts on both sides of the exceptional
curve. Choose the source-base boundary outside the
two moving critical values. Its restriction is then
a proper submersion as well; the boundary transports
used below are transports of this resolved family.
At \(v_0=0\) the I2 has components \(C_0,C_1\);
the marked section meets \(C_1\), and \(O\) meets
\(C_0\). The actual degree-two section closure has
Cartier equation
\[
                  1-Bz^3-(1-2v_0)z^2=0 .
\]
To see this directly, the tangent slope at
\(G=(0,v_0)\) is \(T_0/(2v_0)\), so the actual
double has \(Y+v_0=-T_0X/(2v_0)\). In the first
resolution chart this gives \(T_0=-2v_0/z\).
Substitute \(2v_0=B-Dz\), \(X=Bz\) and
\(D=T_0+X^2+(1-2v_0)X\); canceling on the dense
ordinary chart gives the displayed Cartier equation,
which defines its closure. At \(v_0=0\), one has
\(B=Dz\) and
\(T_0=D(1-z^2-Dz^4)\). The two central components
are therefore \(C_1=\{D=0\}\) and
\(C_0=\{1-z^2-Dz^4=0\}\), with the latter exactly
the displayed double-section closure.
It restricts to \(C_0\) with multiplicity one and
does not contain \(C_1\). Consequently
\[
 c_1(M_1)[C_1]=1,\qquad c_1(M_2)[C_1]=2.               \tag{P3K.57.IE.3}
\]
The second equality retains the actual vertical
Cartier specialization \(2G-O=C_0\).

The local source patch retracts onto its I2
plumbing: use the radial collapse of \(xy=t\)
to the axes in the two node charts and the
component normal-disc retractions elsewhere.
Its integral two-cycle lattice has primitive
basis \(F=C_0+C_1,C_1\), with pairing
\[
                  \begin{pmatrix}0&0\\0&-2\end{pmatrix}.       \tag{P3K.57.IE.4}
\]
The normalized component two-cells give integral
\(C_0,C_1\), and the zero-section relative disc
meets \(F\) once. The smooth source resolution
is a submersion over \(v_0\), also with boundary,
so this marking transports to \(S_{\epsilon_*}\).

On the reference chain (P3K.57.IE.1) the positive normal
line is literally \(M_a\), and the negative one
is \(M_a^{-1}\). Conjugating the negative-weight
line gives determinant
\[
                          b_{\rm r}=2c_1(M_a).          \tag{P3K.57.IE.5}
\]
Indeed the actual fibre-linear change \(V\mapsto gV\)
on \(M_a^\times\) induces
\((X_j,Y_j)\mapsto(gX_j,g^{-1}Y_j)\) in the chain.
After conjugating the negative-weight coordinate,
the determinant transition is
\(g/\overline g=g^2/|g|^2\). A Hermitian metric
identifies the conjugate of \(M_a^{-1}\) with
\(M_a\) by this explicit positive real factor.
The boundary flow \(\Psi\) of IE.1 preserves the
physical circle, so its normal derivative becomes
a complex linear map after this conjugation. Its
full determinant is transported, including any
off-diagonal entries; no diagonal derivative is
assumed.
It evaluates to zero on \(F\) and \(2a\) on \(C_1\).
Its boundary class is zero integrally: in the
central Tate marking \(QP^2=1\), so \(Z^aV^2\)
is a nonvanishing boundary frame. Transport that
actual frame through the smooth source/Cartier
family. Equivalently the boundary component class
has order two and (P3K.57.IE.5) kills it. A relative lift
is \(-aC_1\), up to the null fibre and a boundary
connecting class. Thus
\[
                              b_{\rm r}^2=-2a^2.        \tag{P3K.57.IE.6}
\]

Here is the chart map identifying the collision E
patch with that same minimal I2 patch. The E face
between the primitive rays
\((0,j,1,0)\), \((0,j+1,1,0)\) has quotient lattice
\((Z,s)\), star rays \((0,1),(1,1)\), and primitive
source period \((Z,s)\mapsto(Z+2s,s)\). On the central
source write \(Q=t^2u_Q(t)\) and use the actual
square root with transported leading coefficient,
\(\widetilde q=t\sqrt{u_Q(t)}\). The two-node chart
chain is
\[
 \xi_\ell=Z\widetilde q^{-\ell},\qquad
 \eta_\ell=\widetilde q^{\ell+1}/Z,\qquad
 \xi_\ell\eta_\ell=\widetilde q,
 \quad \ell\pmod2.                                  \tag{P3K.57.IE.5a}
\]
Its literal overlaps are
\(\xi_{\ell+1}=\eta_\ell^{-1}\),
\(\eta_{\ell+1}=\xi_\ell\eta_\ell^2\).
These are the E-star characters with the retained
Tate units; the primitive stabilizer shifts
\(\ell\) by two. In particular the map has no
additional base ramification or exceptional curve.

For completeness, compose this map with the actual
Tate uniformization into the unit source
\(Y_0^2=X_0^3+X_0^2+T_0X_0\). In its marked Tate
coordinate the convergent functions are
\[
 X_T=\sum_{\nu\in\mathbb Z}
       \frac{Q^\nu Z}{(1-Q^\nu Z)^2}
       -2\sum_{\nu\ge1}\frac{\nu Q^\nu}{1-Q^\nu},
 \qquad
 Y_T=\sum_{\nu\in\mathbb Z}
       \frac{(Q^\nu Z)^2}{(1-Q^\nu Z)^3}
       +\sum_{\nu\ge1}\frac{\nu Q^\nu}{1-Q^\nu}.
\]
The retained minimal Weierstrass coordinate change
has leading terms
\(X_0=4X_T-T_0/2+O(T_0^2,T_0X_T)\) and
\(Y_0=8Y_T+4X_T+O(T_0Y_T,T_0X_T)\), while
\(T_0=16\widetilde q+O(\widetilde q^2)\).
On the chart \(Z=\xi\),
\(\widetilde q=\xi\eta\), this gives
\[
 T_0=16\xi\eta+O(\xi^2),\quad
 X_0=4\xi(1-\eta)^2+O(\xi^2),\quad
 Y_0=4\xi(1-\eta^2)+O(\xi^2).                       \tag{P3K.57.IE.5b}
\]
For every negative summation index, rewrite the
summand in \(Q^k/Z\); its exponents are nonnegative
in this node chart. The resulting series converge
normally on a smaller chart, so these are actual
holomorphic maps, including the axes.

The map to the two small-resolution charts is
\[
 z=X_0/Y_0,\quad D=T_0+X_0+X_0^2,\quad B=Y_0,
 \qquad
 w=Y_0/X_0,\quad C=X_0,\quad A=Y_0.                 \tag{P3K.57.IE.5c}
\]
The source equation gives \(A=Dz,C=Bz\) exactly.
At the marked two-torsion point \(\eta=1\), the
double zero of \(X_0\) and simple zero of \(Y_0\)
make the first quotient extend. On \(C_1\),
\[
 z=\frac{1-\eta}{1+\eta},\qquad
 D=4\xi(1+\eta)^2+O(\xi^2).
\]
The Jacobian of \((\xi,\eta)\mapsto(D,z)\) on
\(\xi=0\) is \(-8\). At \(\eta=-1\), use
\((C,w)\): its Jacobian has magnitude eight there.
The second node is covered by the literal chain
overlap in (P3K.57.IE.5a). This proves the biholomorphism
on the whole component, its two nodes, and the
marked section, with the source orientation and
boundary marking retained.

The actual normal determinant is marked by
\[
                         Z^aV^2/\xi_E^{\,2j+1}.
\]
After extracting its weight-two normal factor,
this character is a meromorphic section of
\(N_+\otimes N_-^{-1}\). It is invariant under the
source period, since \(QP^2=1\), and is nonzero
away from the two central components. Its
dual-character order is zero on each component.
The actual node charts above therefore make the
section invertible also at their intersections.
The Hermitian identification of this line with
the positive-weight spinor determinant is the
same one used in (P3K.57.IE.5). This proves zero
integral first Chern class, rather than merely
a rational vanishing of its degree. One can also
read the integral obstruction from the plumbing
retraction: \(H_1=\mathbb Z\), \(H_2=\mathbb Z^2\)
with primitive basis \(C_0,C_1\), so
\(\operatorname{Ext}(H_1,\mathbb Z)=0\); vanishing
on both components determines the zero class
in \(H^2\). Hence in this boundary marking
\[
                               b_{\rm c}^2=0.           \tag{P3K.57.IE.7}
\]
Replacing \(\xi_E=\epsilon/t\) by \(\epsilon\)
adds the explicit boundary gauge \(t^{2j+1}\).
It may be odd and is retained. A boundary
gauge changes the glued class by the
Mayer--Vietoris image of \(H^1(\partial S;\mathbb Z)\).
Its square and its product with the glued
determinant class vanish. Indeed represent
the connecting class as \(d\chi\wedge\alpha\)
in the collar; its square repeats \(d\chi\),
and its product integrates the zero boundary
restriction of the other class. Thus the
actual derivative of (P3K.57.IE.2) has exactly the
square difference (P3K.57.IE.7) minus (P3K.57.IE.6).

The source patches are oriented diffeomorphic
with boundary. Their closed double has
signature zero and hence tangent Pontryagin
integral zero. This uses only the ordinary
four-dimensional fixed-surface signature
identity. With weights \(+1,-1\), its
normal-link contribution is
\[
 \frac1{24}\int_{S_{\rm c}\cup(-S_{\rm r})}
            \{p_1(TS)-3b^2\}=-\frac{a^2}{4}.           \tag{P3K.57.IE.8}
\]
There are exactly \(N\) E doubles: the actual
scalar period shifts the primitive chain by
\(N\), and the collision-boundary source
loop fixes these labels.

\par\medskip\noindent\textbf{IE.4. Finite stabilizers and the actual Stokes identity.}\par\nopagebreak[4]

Remove invariant normal tubes of every fixed
component from \(W_{\rm rel}\). The physical
circle is locally free on the complement.
Its quotient is a compact oriented
seven-dimensional orbifold \(Z\), carrying
the actual orbifold Euler line with class \(x\).

A slice with stabilizer \(\mu_d\) gives a
finite chart \(\widetilde U/\mu_d\). Average
metrics, partitions of unity and connections
over \(\mu_d\) and over the physical circle.
Normalize the connection by
\(\alpha(\delta)=1\). Its curvature represents
the same physical \(x\) on every chart. The
actual tangent splitting is
\[
 TW_{\rm rel}|_{\rm locally\ free}
                    \cong\pi^*TZ\oplus\mathbb R_\delta .      \tag{P3K.57.IE.9}
\]
Choose its connection to match the actual
boundary tangent connections after a collar
interpolation. The basic Chern--Weil form
\[
                      \Omega=\frac{4x^3-p_1(TZ)x}{24}         \tag{P3K.57.IE.10}
\]
is closed. Integrate each finite chart as
\(d^{-1}\int_{\widetilde U}\). Ordinary Stokes
on the finite covers and their averaged
partition of unity cancels every interior
face with exactly these multiplicities.
The finite-isotropy strata are not boundary
faces. Consequently
\(\int_{\partial Z}\Omega=0\).

On a fixed normal link, (P3K.57.IE.9) restricts
to the fixed tangent and weighted normal
bundles. Its equivariant Thom denominator
is \(\prod_i(x_i+w_i u)\). The finite
chart integral, equivalently the weighted
projective power map, gives
\(1/\prod_iw_i\), with every negative
normal conjugation's orientation sign.
Thus its actual contribution is (P3K.57.IE.16), proved below.
The compact interior/H values and (P3K.57.IE.8)
are exactly the terms in this Stokes
identity, including finite stabilizers.
No spin structure on this comparison
bordism or eight-dimensional signature
formula is used.

\par\medskip\noindent\textbf{IE.5. Outer boundary and Jin conventions.}\par\nopagebreak[4]

Orient \(W_{\rm c}\) complexly and the
reference piece oppositely. Define
\(\delta\wedge\operatorname{or}(Z)
=\operatorname{or}(W_{\rm rel})\).
Near the outer parameter boundary,
\[
 \operatorname{or}(W)
 =\delta\wedge\operatorname{or}(Y)
                  \wedge d\rho\wedge d\theta .
\]
Since \(\dim Y=5\), its quotient boundary
orientation is
\[
                         -\operatorname{or}(Y)\wedge d\theta.       \tag{P3K.57.IE.11}
\]
At an inner fixed-normal boundary the
minus sign from removing the normal disc
and the minus sign from quotienting an
oriented boundary by the odd-dimensional
orbit cancel. Its orientation is the
positive complex-normal link convention
of (P3K.57.IE.16). Therefore Stokes gives
\[
                 A_{\rm physical,difference}
                          =\sum_F\mathcal R(F).          \tag{P3K.57.IE.12}
\]

Jin identifies \((z,0)\sim(f(z),1)\).
The positive physical family with return
\(f\) has the inverse parameter marking.
The map to the Jin six-manifold reverses
that circle and preserves the physical
Euler line. Hence
\[
                  A_{\rm Jin,difference}
                           =-\sum_F\mathcal R(F).         \tag{P3K.57.IE.13}
\]
The common conifold replacement in IE.2
uses the same transported positive
meridian, so adds no line or orientation
change.

\par\medskip\noindent\textbf{IE.6. Scalar-order independence and exact local values.}\par\nopagebreak[4]

The selected-cusp fan in this argument is the initial
vertical-edge fan or its literal reflected fan under
\((Z,V)\mapsto(Z^{-1},P^aV)\). In either marking the
physical delta circle has the already specified
\((-1,+1)\) fixed-curve model, and the selected-unit
transport in IE.2 closes that same fan. The admitted
source/cap sectors are \((a,b)=(1,2)\), \(m=4p\),
and \((a,b)=(2,1)\), \(m=4p-1\), with \(p\ge1\),
their actual \(n=\pm1\) norm markings, and the
retained integral \(d\delta\) twist. No comparison
with a further flopped fan is used here.

The fixed-source scalar integers in these two particular markings
are equal by the actual C2 comparison of P3K.62. Its first outer
branch map is the identity; its second has base multiplier
\(\kappa^{-1}\) and normal matrix
\(\operatorname{diag}(\alpha^{-1},\kappa\alpha)\), both independent
of the scalar parameter. These are exactly the outer maps used
on the selected collar here. Together with the equal inner
charged-cap gauges, they give the line-preserving integer
comparison in P3K.62. This invocation retains the two outer
maps and uses no assertion of equality for a further fan.

Over \(U\) the ordinary suspension, with
no linearization zero, is an actual family
over the product of a source annulus and
a small scalar annulus:
\[
                   B_\lambda=2(-1)^{b+1}\lambda q_*U_a^b.
\]
There is no finite cap in this local
family, so every integer scalar winding
is a closed local loop. Its source-base
boundary family extends over the source
disc, with the actual periods and Cartier
frames. Choose IE.1's transport on this
entire product.

A pair of pants mapping into this
product realizes the based-loop identity
\[
 (\epsilon\text{ turn},N\text{ scalar turns})
   =(\epsilon\text{ turn},0)+(0,N\text{ scalar turns}).    \tag{P3K.57.IE.14}
\]
Pull back the actual local family and
quotient its free physical circle.
This is a seven-dimensional bordism
of the mapping-torus pieces. Glue the
source-base boundary collars by the
full-product transport. The fixed-source
scalar-loop boundary is exactly the
reference piece of IE.1 and cancels.
Thus the closed relative collision
boundaries \(Q_{\rm rel}(N)\) and
\(Q_{\rm rel}(0)\) are bordant with
their actual Euler line. Their
characteristic difference is independent
of \(N\). This is a literal parameter
product and collar construction; no
division of a covered sphere class occurs.

For any admitted cap, its exterior
source/root/line family extends over
the source disc. IE.2 cancels it with
the constant exterior. Hence this
same local relative boundary represents
the fixed-scalar source loop for that
completed family. It suffices to use
the two fully executed orders \(N=8\)
for \((a,b)=(1,2)\), and \(N=6\) for
\((a,b)=(2,1)\):
\[
 \begin{array}{c|rrr}
 (a,b,N)&\text{interior/H}&N\text{ E doubles}&\text{total}\\ \hline
 (1,2,8)&2&-2&0\\
 (2,1,6)&5&-6&-1 .
 \end{array}                                           \tag{P3K.57.IE.15}
\]
Independent executions at \(N=4,12\)
in the first case and \(N=3,9\) in
the second give the same relative
sums. They check the arithmetic;
the independence proof is (P3K.57.IE.14).

By (P3K.57.IE.13), the fixed-scalar marked infinity
Jin differences are respectively \(0\) and \(+1\).
Apply P3F.22's proved Jin-to-\(\Theta_7\) transfer
to the two actual completed boundary mapping
tori: on each, IE.2 has made the specified conifold
replacement, transported the same positive meridian,
and retained the actual Euler line of its free circle
bundle. To select their actual return maps, prescribe
the product connection on the common selected tube
in IE.2's closed Cox product, extend it to a connection
on the actual parameter family, and average over
the physical delta circle. The prescribed product
connection remains unchanged under that averaging.
After the conifold replacement its return is a
diffeomorphism \(f\) of \(Y=S^2\times S^3\), equal
to the identity on the prescribed tube. The actual
Euler line extends around the family, so \(f\)
preserves its primitive class \(x\in H^2(Y;\mathbb Z)\).
The return preserves orientation; the perfect
pairing \(H^2\otimes H^3\to H^5\) then shows that
it preserves \(H^3\) as well. Thus \(f\in I(Y)\).

On the complement of that tube, P3F.22's Hopf lift
is literally the physical return lift supplied by
this delta-equivariant connection and its preserved
Euler line. On the tube it is the same identity
extension in the proved conifold/Hopf collar
construction. Consequently its \(\lambda(f)\) is
the actual clutched sphere of this parameter return.
This identifies the return map and its lift before
using the characteristic polynomial. Their
characteristic integers are
\(\int(4x^3-p_1x)/24\) in the Jin convention.
The transfer sends their difference to that integer
modulo 28. The pair-of-pants map (P3K.57.IE.14) identifies
the difference of these actual returns with the
fixed-scalar positive epsilon source loop. Thus
its oriented sphere class is \(0\) for \((a,b)=(1,2)\)
and \(1\) for \((a,b)=(2,1)\), in each of the two
specified vertical-edge fan markings. This conclusion
does not assign either source-zero fan return.

The old
\(v\)/new \(\epsilon\) scalar path,
\(D^{-4}\) root endpoint and
half-loop fan map remain the explicit
dictionary of IF.1--6. Their winding
corrections are not absorbed into
(P3K.57.IE.15).

\par\medskip\noindent\textbf{IE.7. Complete normal-link calculation for the two executed fans.}\par\nopagebreak[4]

This section supplies the local formulas used in IE.4 and
the exact finite sums in (P3K.57.IE.15), independently of the
uncompleted source-zero attachment. Let a compact fixed
component \(F\) have complex normal roots \(x_i\) and
nonzero physical delta weights \(w_i\). Retain its complex
orientation and the original complex normal-sphere
orientation. Its actual normal-link quotient contribution is
\[
 \mathcal R(F)=[u^{-1}]\int_F
 \frac{4u^3-u\{p_1(TF)+\sum_i(x_i+w_i u)^2\}}
      {24\prod_i(x_i+w_i u)}.                         \tag{P3K.57.IE.16}
\]
The coefficient is the Laurent coefficient at infinity.
On the normal sphere, adding its radial real line gives
the pullback of \(TF\oplus\nu_F\); quotienting the
physical orbit removes its oriented real line. Thus the
stable quotient Pontryagin class is the displayed
equivariant class. The equivariant Thom class restricts
to \(\prod_i(x_i+w_i u)\). The localized pushforward
therefore gives (P3K.57.IE.16), with the Euler and integration
identities of Atiyah--Bott \cite[pp.~8--9, Theorem~(3.5), equations~(3.7)--(3.8)]{AtiyahBott1984}.

The fibre normalization can be verified directly. For
positive weights, the weighted projective power map
\([z_i]\mapsto[z_i^{w_i}]\) has generic degree
\(\prod_iw_i\) and pulls its rational hyperplane class
back to the ordinary hyperplane class. Consequently
\(\int u^{r-1}=1/\prod_iw_i\). The relation
\(\prod_i(x_i+w_i u)=0\) gives all higher pushforwards
recursively. Conjugating a negative-weight line changes
\((x_i,w_i)\) to \((-x_i,-w_i)\) and reverses the
original normal orientation. This is precisely the
denominator sign in (P3K.57.IE.16); every finite stabilizer
factor remains.

For a point and a curve the respective formulas are
\[
 \frac{4-\sum_iw_i^2}{24\prod_iw_i},\qquad
 -\frac1{24\prod_iw_i}\int_F\sum_i
 \left\{\frac{4-\sum_jw_j^2}{w_i}+2w_i\right\}x_i.
                                                               \tag{P3K.57.IE.17}
\]
For a surface, put
\[
 P=w_1w_2,\quad c=4-w_1^2-w_2^2,\quad
 A=\frac{x_1}{w_1}+\frac{x_2}{w_2},\quad
 B=\frac{x_1^2}{w_1^2}+\frac{x_1x_2}{w_1w_2}
                         +\frac{x_2^2}{w_2^2}.
\]
Its contribution is
\[
 \int_F\frac{cB+2(w_1x_1+w_2x_2)A
                 -p_1(TF)-x_1^2-x_2^2}{24P}.             \tag{P3K.57.IE.18}
\]
Weights \((1,-1)\) give
\((p_1(TF)-3(x_1-x_2)^2)/24\), the formula used
in (P3K.57.IE.8), in the same original normal orientation.

For each exported regular cone let \(C\) be the
ordered matrix of its four primitive rays in the
original \((Z,V,\text{base}_0,\text{base}_1)\)
lattice. The coordinate characters are the rows of
\(C^{-1}\), and their actual physical weights are
\[
                    C^{-1}(0,1,0,0)^{\mathsf T}.        \tag{P3K.57.IE.19}
\]
The vanishing of the nonzero-weight coordinates
specifies every local fixed stratum. The recorded
face permutations and literal period matrices glue
these strata, so no component is counted twice or
omitted. There is no fixed stratum in the open
protected A or B face: their span is respectively
\(\langle(1,0,0,0),(0,0,1,0)\rangle\) or
\(\langle(1,-a,0,0),(0,0,1,0)\rangle\), neither
of which contains the physical delta vector.
Their remaining boundary strata belong to the
already recorded H/interior cones. The protected
E face does contain delta; its fixed surfaces are
exactly the N chain components treated in IE.3.
Thus the interior, H and E contributions exhaust
the fixed set of the relative comparison.
On an interior fixed component, unwrap
each chart by its recorded period \(n\). An
auxiliary cocharacter \(e\) then has tangent and
normal weights
\[
                            C^{-1}T_n e.                 \tag{P3K.57.IE.20}
\]
This is the dual of the actual map
\(C_{\rm unwrapped}=T_{-n}C\). Applying the same
pushforward to the fixed component divides (P3K.57.IE.17)
or (P3K.57.IE.18), evaluated at its auxiliary roots, by
the product of its tangent weights. The two
independently executed choices are
\((1,4,16,64)\) and \((1,5,25,125)\); the code
verifies agreement separately on every component.

On the exceptional base H use \(w=q_B/q_A\) and
normal coordinate \(q_A\). The exact periods have
normal-order matrix
\[
 H_\infty=\begin{pmatrix}2&b\\-a&N-2\end{pmatrix},
\]
and remaining nonconstant factors
\[
 h:(Z,V)\mapsto(wZ,w^{-a}V),\qquad
 g:(Z,V)\mapsto(w^bZ,w^{-2}(1-w)^N V).                 \tag{P3K.57.IE.21}
\]
All other factors are the retained nonzero source,
Cartier and suspension units of IF.5. The three
marked base points are \(w=0,1,\infty\), on the
B,E,A faces respectively. For each regular rank-three
H face, drop its zero first base coordinate and
form the unimodular matrix \(B_H\). Its orbit
closure has the actual base coordinate \(w\), so
is \(\mathbb P^1\). If
\((p_j,q_j,s_j)\) is row \(j\) of \(B_H^{-1}\),
its generic normal coordinate is
\(\eta_j=Z^{p_j}V^{q_j}q_A^{s_j}\).

For an occurrence at period \(n=(n_1,n_2)\), let
\((p,q,r,s)\) be its corresponding full chart row.
The exact recorded checks are
\((p,q)=(p_j,q_j)\) and
\(s+(p,q)H_\infty n=s_j\). Set
\[
 e_w=p(n_1+bn_2)+q(-an_1-2n_2),\qquad e_1=Nn_2q.
\]
The three actual coefficient orders of the generic
normal vector, in its holomorphic local frames, are
\[
 \begin{array}{c|cc|c}
  \text{sector}&x&y&\text{order at its marked point}\\ \hline
  B&w&q_A&e_w+r\\
  E&(1-w)/u_\epsilon&q_A&e_1+r\\
  A&w^{-1}&wq_A&-e_w+r-s-e_1.
 \end{array}                                             \tag{P3K.57.IE.22}
\]
Their sum is the actual normal-line degree. In
particular the \((1-w)^N\) term, the higher ray
heights, and every recorded period remain in the
degree computation.

The two auxiliary fibre circles \((1,2,0,0)\) and
\((1,3,0,0)\) commute with (P3K.57.IE.21) and the full
periods. On each compact H fixed surface, their
fixed sets consist of these rank-three curves and
the isolated chart fixed points. A sector-interior
three-face has base rank two and only the physical
one-dimensional fibre kernel, so supplies no
additional auxiliary fixed curve. If an auxiliary
fixed curve has normal weight \(\eta\), normal
degree \(d_t\), and ambient normal weights/degrees
\((\beta_i,d_i)\), substitute
\[
 x_i=\beta_i u+d_i h,\qquad
 p_1(TF)=\eta^2u^2+2\eta d_t hu,\qquad h^2=0
\]
into (P3K.57.IE.18). Writing the numerator polynomial
there as \(R_0u^2+R_hhu\), its exact localized
contribution is
\[
                         R_h/\eta-R_0d_t/\eta^2.         \tag{P3K.57.IE.23}
\]
This divides by the actual Euler class
\(\eta u+d_t h\). Isolated fixed points use the
two nonzero tangent weights. Both auxiliary circles
agree component by component. Independent outputs
also compute \(c_1(TF)^2,c_2(TF)\), and verify the
integrality of \((c_1^2-2c_2)/3\) and
\((c_1^2+c_2)/12\) on every H surface.

The complete finite outputs are
\[
 \begin{array}{c|rr|rr|r}
 (a,b,N)&\#\text{ interior}&\sum_{\rm int}\mathcal R
 &\#H&\sum_H\mathcal R&\text{combined}\\ \hline
 (1,2,8)&2153&9333/16&125&-9301/16&2\\
 (2,1,6)&2147&739&111&-734&5.
 \end{array}                                             \tag{P3K.57.IE.24}
\]
There are 194 and 152 full-fibre auxiliary fixed
curves on H, respectively. Every primitive ray,
ordered determinant-one cone, period, face
permutation, normal degree and rational summand
is retained in the corresponding
\(\texttt{source\_infinity\_regular\_fan}\),
\(\texttt{fixed\_charts}\),
\(\texttt{interior\_residues}\), and
\(\texttt{h\_residues}\) JSON/certificate files
under \(\texttt{code/part3\_fourfold/}\).
Their scope is the actual infinity matrices and
period maps just displayed. Combining (P3K.57.IE.24)
with the E doubles proved in IE.3 gives (P3K.57.IE.15).
\end{proof}
